%──────────────────────────────────────────────────────────────
%  CV — Daniël van Ginneken (Nederlands)
%──────────────────────────────────────────────────────────────
\documentclass[a4paper,10pt]{article}
%──────────────────────────────────────────────────────────────
%  Shared CV preamble — Daniël van Ginneken
%  \input this from each language's main .tex file.
%  Keep packages, colours, section/entry formatting here so the
%  EN and NL CVs can never drift out of visual sync.
%
%  Layout: a dark header band, then a two-column grid — a narrow
%  "meta" column (dates, locations, labels) on the left and the
%  content column on the right. Every section uses the same grid.
%  Single text flow, no floats/columns, so ATS parsers read it
%  top-to-bottom.
%
%  Engine: pdflatex (same as CI). Not XeTeX: there, Source Sans hyphens
%  extract from the PDF as U+2011, which breaks ATS keyword matching,
%  and microtype can't do font expansion.
%  Single pass — no tikz overlays that need a previous run's .aux.
%──────────────────────────────────────────────────────────────

% ── Packages ─────────────────────────────────────────────────
\usepackage[a4paper, left=1.2cm, right=1.2cm, top=0.7cm, bottom=0.6cm, noheadfoot]{geometry}
\usepackage[T1]{fontenc}
\usepackage[utf8]{inputenc}
\usepackage[default, semibold]{sourcesanspro}
\usepackage{titlesec}
\usepackage{enumitem}
\usepackage{xcolor}
\usepackage{fontawesome5}
\usepackage{microtype}
\usepackage{graphicx}
\usepackage{tikz}
\usepackage{eso-pic}
\usepackage{hyperref}

% Letterspacing (tracking) via microtype; #1 is in 1/1000 em.
\newcommand{\track}[2]{\textls[#1]{#2}}

% ── Colours ──────────────────────────────────────────────────
\definecolor{band}{HTML}{16222D}      % header background
\definecolor{bandtext}{HTML}{D5DEE5}  % contact text on the band
\definecolor{accent}{HTML}{1C7C8C}    % teal: headings, links, markers
\definecolor{accentlight}{HTML}{6CC4D1}
\definecolor{ink}{HTML}{1E2A33}       % body text
\definecolor{muted}{HTML}{5E6B75}     % dates, secondary text
\definecolor{ruleclr}{HTML}{D3DAE0}

\color{ink}

% ── Hyperlinks ───────────────────────────────────────────────
\hypersetup{
  colorlinks=true,
  urlcolor=accent,
  linkcolor=accent,
  pdfauthor={Daniël van Ginneken},
  pdftitle={CV — Daniël van Ginneken},
}

% ── Page basics ──────────────────────────────────────────────
\pagestyle{empty}
\setlength{\parindent}{0pt}
\setlength{\parskip}{0pt}
\AtBeginDocument{\fontsize{9.6}{11.6}\selectfont}
\hyphenpenalty=5000
\exhyphenpenalty=5000

% ── Header band ──────────────────────────────────────────────
% Painted behind page 1 via eso-pic; the header text sits on top of
% it in the normal flow, inside a box of fixed height.
\newlength{\bandheight}\setlength{\bandheight}{3.45cm}
\newlength{\bandpad}\setlength{\bandpad}{0.7cm}   % = geometry top

\AddToShipoutPictureBG*{%
  \AtPageUpperLeft{%
    {\color{band}\rule[-\bandheight]{\paperwidth}{\bandheight}}%
  }%
  \AtPageUpperLeft{%
    \raisebox{-\bandheight}{{\color{accent}\rule[-3pt]{\paperwidth}{3pt}}}%
  }%
}

% \contactlink{icon}{url}{text} — light-on-dark contact item
\newcommand{\contactlink}[3]{%
  \mbox{{\color{accentlight}#1}\,\href{#2}{\color{bandtext}#3}}}
\newcommand{\contactplain}[2]{%
  \mbox{{\color{accentlight}#1}\,{\color{bandtext}#2}}}
\newcommand{\csep}{\hspace{0.7em}}

% \cvheader{name}{headline}{contact line 1}{contact line 2}{photo}
\newcommand{\cvheader}[5]{%
  \begin{minipage}[c][\dimexpr\bandheight-2\bandpad\relax][c]{\dimexpr\textwidth-3.3cm\relax}
    {\color{white}\fontsize{27}{30}\selectfont\bfseries #1}\par\vspace{4pt}
    {\color{accentlight}\large\track{40}{#2}}\par\vspace{9pt}
    {\small #3}\par\vspace{3pt}
    {\small #4}
  \end{minipage}%
  \hfill
  \begin{minipage}[c][\dimexpr\bandheight-2\bandpad\relax][c]{3cm}
    \raggedleft
    \begin{tikzpicture}
      \begin{scope}
        \clip circle (1.32cm);
        \node[yshift=-0.15cm] {\includegraphics[width=2.75cm]{#5}};
      \end{scope}
      \draw[accentlight, line width=1.6pt] circle (1.32cm);
    \end{tikzpicture}
  \end{minipage}\par
  \vspace{\dimexpr\bandpad+8pt\relax}%
}

% ── Section headings ─────────────────────────────────────────
% TITLE
% ━━━━━━━━━━━━━━━━━━━━━━━━━━━━━━━━━━━━━━ (solid dark rule, full width)
\newcommand{\sectionlabel}[1]{\track{60}{\MakeUppercase{#1}}}
\titleformat{\section}[block]
  {\color{band}\sffamily\bfseries\large}
  {}{0pt}{\sectionlabel}[\vspace{-1pt}{\color{band}\titlerule[0.9pt]}]
\titlespacing*{\section}{0pt}{7pt}{4.5pt}

% ── The grid ─────────────────────────────────────────────────
\newlength{\metacol}\setlength{\metacol}{2.65cm}
\newlength{\metagap}\setlength{\metagap}{0.3cm}

% Every section body is one `cvgrid`; each \item[meta] is a row.
\newlist{cvgrid}{description}{1}
\setlist[cvgrid]{
  leftmargin=\dimexpr\metacol+\metagap\relax,
  labelwidth=\metacol, labelsep=\metagap, labelindent=0pt,
  align=left, font=\normalfont,
  topsep=0pt, itemsep=4.5pt, parsep=0pt, partopsep=0pt,
  before=\raggedright,
}
% Two-line meta cell: primary (dates) + secondary (location)
\newcommand{\meta}[2]{%
  \smash{\parbox[t]{\metacol}{\raggedright\leavevmode
    {\small\color{muted}#1}%
    \if\relax\detokenize{#2}\relax\else\\[-1pt]{\footnotesize\color{muted}#2}\fi}}}

% \cventry{dates}{location}{role}{organisation}
\newcommand{\cventry}[4]{%
  \item[\meta{#1}{#2}]%
  {\bfseries #3}\\[-1pt]
  {\color{accent}#4}\par}

% ── Several roles at one employer ───────────────────────────
%   \cvcompany{total dates}{location}{organisation}
%   \cvrole{role}{role dates}{bullets}     — newest first
%   \cvrole*{role}{role dates}{bullets}    — the last (oldest) role
%
% Three looks for the same content. Pick one by (un)commenting the
% \NewCommandCopy lines at the end of this block — the .tex files
% don't change.

% (1) GROUPED — company once (total span in the meta column), roles as
%     sub-headings with their own dates right-aligned.
\newcommand{\groupedcompany}[3]{%
  \item[\meta{#1}{#2}]%
  {\color{accent}#3}\par\vspace{1pt}}
\NewDocumentCommand{\groupedrole}{s m m +m}{%
  {\bfseries #2}\hfill{\small\color{muted}#3}\par
  #4%
  \IfBooleanF{#1}{\par\vspace{4pt}}}

% (2) TIMELINE — like (1), plus a LinkedIn-style dot per role and a rail
%     connecting them. The rail is a \vrule beside a [t] minipage, so
%     it's exactly as tall as the role's content.
\newlength{\railgap}\setlength{\railgap}{9pt}
\newcommand{\raildot}{%
  \llap{\tikz[baseline=-0.55ex]\fill[accent] (0,0) circle (2.3pt);\hspace{7.4pt}}}
\NewDocumentCommand{\timelinerole}{s m m +m}{%
  \par\nointerlineskip\noindent
  \hbox{%
    \hspace{3pt}%
    \IfBooleanTF{#1}{\hspace{0.8pt}}{{\color{ruleclr}\vrule width 0.8pt}}%
    \hspace{\railgap}%
    \begin{minipage}[t]{\dimexpr\linewidth-3.8pt-\railgap\relax}
      \raggedright
      \leavevmode\raildot{\bfseries #2}%
      \hfill{\small\color{muted}#3}\par
      #4%
      % trailing gap kept inside the box so the rail bridges it
      \IfBooleanF{#1}{\par\vspace{4pt}\hrule height 0pt}%
    \end{minipage}}%
  \par}

% (3) SEPARATE — every role is its own entry with the company repeated,
%     role dates in the meta column (the classic layout).
\newcommand{\separateorg}{}\newcommand{\separateloc}{}
\newcommand{\separatecompany}[3]{%
  \renewcommand{\separateloc}{#2}\renewcommand{\separateorg}{#3}}
\NewDocumentCommand{\separaterole}{s m m +m}{%
  \cventry{#3}{\separateloc}{#2}{\separateorg}#4}

% ▸ Active look — keep exactly one pair uncommented:
\NewCommandCopy\cvcompany\groupedcompany  \NewCommandCopy\cvrole\groupedrole     % (1) grouped
% \NewCommandCopy\cvcompany\groupedcompany  \NewCommandCopy\cvrole\timelinerole  % (2) timeline
% \NewCommandCopy\cvcompany\separatecompany \NewCommandCopy\cvrole\separaterole  % (3) separate

% \cvproject{year}{name}{url}{display url}{stack}
\newcommand{\cvproject}[5]{%
  \item[\meta{#1}{}]%
  {\bfseries #2}\hspace{0.5em}{\small\href{#3}{#4}}%
  \if\relax\detokenize{#5}\relax\else\\[-1pt]{\small\color{muted}#5}\fi\par}

% \cvrow{label}{content} — skills, languages, certifications
\newcommand{\cvrow}[2]{%
  \item[{\small\bfseries #1}]#2}

% Interpunct separator for inline lists
\newcommand{\middot}{\unskip\hspace{0.4em}{\color{accent}\textperiodcentered}\hspace{0.4em}\ignorespaces}

% ── Bullets ──────────────────────────────────────────────────
\setlist[itemize]{
  label={\color{accent}\raisebox{0.18em}{\rule{0.3em}{0.3em}}},
  leftmargin=11pt, labelsep=6pt,
  topsep=2pt, itemsep=0.5pt, parsep=0pt, partopsep=0pt,
}


\begin{document}

% ═══════════════════════════════════════════════════════════
% HEADER
% ═══════════════════════════════════════════════════════════
\cvheader{Daniël van Ginneken}
  {Backend Softwareontwikkelaar}
  {\contactplain{\faMapMarker*}{Oudenbosch, Nederland}\csep
   \contactlink{\faEnvelope}{mailto:dcjvanginneken@gmail.com}{dcjvanginneken@gmail.com}\csep
   \contactplain{\faPhone*}{Beschikbaar op aanvraag}}
  {\contactlink{\faLinkedin}{https://linkedin.com/in/danielvanginneken}{linkedin.com/in/danielvanginneken}\csep
   \contactlink{\faGlobe}{https://danielvanginneken.nl}{danielvanginneken.nl}\csep
   \contactlink{\faGithub}{https://github.com/danielvanginneken}{github.com/danielvanginneken}}
  {../assets/picture.jpg}

% ═══════════════════════════════════════════════════════════
% PROFIEL
% ═══════════════════════════════════════════════════════════
\section{Profiel}

Backend-gerichte softwareontwikkelaar met 3+ jaar productie-infrastructuurervaring en een sterke interesse in gedistribueerde systemen, event-driven architecturen en backend API-ontwerp.
Thuis in de volledige deployment stack --- van Cisco/VLAN-netwerken en Linux-serverbeheer tot ASP.NET Core API's en message-driven backends.
Studeert momenteel Informatica (HBO) aan Avans Hogeschool en loopt stage als developer bij Basic-Fit.

% ═══════════════════════════════════════════════════════════
% WERKERVARING
% ═══════════════════════════════════════════════════════════
\section{Werkervaring}

\begin{cvgrid}
\cventry{aug 2026 -- jan 2027}{Tilburg, NL}
        {Fullstack Developer \textnormal{\color{muted}$\cdot$ Stage}}{Basic-Fit}
\begin{itemize}
  \item Bouw fullstack features (TypeScript) voor een klantgerichte retailapplicatie, met integratie van externe partijen
  \item Werk binnen een platformteam mee aan refinements en beveiligingsvraagstukken
  \item Documenteer het applicatielandschap en kernprocessen voor het team
\end{itemize}

\cvcompany{nov 2022 -- aug 2026}{Etten-Leur, NL}{Dentech Telecommunicatie, Beveiliging \& Automatisering}
\cvrole{System Support Engineer \textnormal{\color{muted}$\cdot$ Parttime}}{mei 2024 -- aug 2026}{%
\begin{itemize}
  \item Ontwerp en beheer van multi-site netwerkinfrastructuur: Cisco-switching, VLAN-segmentatie en firewallbeleid
  \item Beheer van Linux- en Windows Server-omgevingen (AD, Group Policy, DNS/DHCP) voor diverse mkb-klanten
  \item Implementatie en beheer van Proxmox-virtualisatieclusters met productie-workloads
  \item Automatisering van provisioning- en configuratietaken met PowerShell en Bash, waardoor repetitieve handmatige stappen worden geëlimineerd
  \item Leiding over implementatie en configuratie van 3CX PBX-systemen en VoIP-infrastructuur bij klanten
\end{itemize}}
\cvrole*{Junior Support Engineer \textnormal{\color{muted}$\cdot$ BBL}}{nov 2022 -- mei 2024}{%
\begin{itemize}
  \item Oplossen van hardware-, software- en VoIP-incidenten in meerdere klantomgevingen
  \item Ondersteuning van Windows Server, Active Directory en werkstationbeheer onder begeleiding van seniors
  \item Kerninfrastructuurvaardigheden opgebouwd tijdens MBO Niveau 4 ICT Systeem- \& Devicetechnicus
\end{itemize}}
\end{cvgrid}

% ═══════════════════════════════════════════════════════════
% PROJECTEN
% ═══════════════════════════════════════════════════════════
\section{Projecten}

\begin{cvgrid}
\cvproject{2026}{PatientPingeling}
          {https://github.com/PatientPingeling/PatientPingeling}{github.com/PatientPingeling/PatientPingeling}
          {ASP.NET Core \middot PostgreSQL \middot RabbitMQ \middot OpenTelemetry \middot Docker}
\begin{itemize}
  \item Multi-tenant notificatieplatform met per tenant geconfigureerde e-mail-, SMS- en pushkanalen
  \item Clean architecture met een provider-abstractielaag (factory pattern), zodat notificatiekanalen inplugbaar zijn
  \item Asynchrone berichtverwerking via RabbitMQ, waarmee notificatieverzending wordt losgekoppeld van de request lifecycle
  \item Distributed tracing, metrics en gestructureerde logging via OpenTelemetry; deployment met Docker Compose
\end{itemize}

\cvproject{2025}{Persoonlijke Portfolio}
          {https://danielvanginneken.nl}{danielvanginneken.nl}
          {}
\begin{itemize}
  \item Persoonlijke website gebouwd met React en TypeScript; self-hosted via Docker op eigen infrastructuur
\end{itemize}
\end{cvgrid}

% ═══════════════════════════════════════════════════════════
% OPLEIDING
% ═══════════════════════════════════════════════════════════
\section{Opleiding}

\begin{cvgrid}
\cventry{sep 2024 -- jun 2028}{Breda, NL}
        {Bachelor of Applied Science -- Informatica (HBO)}{Avans Hogeschool}
{\small\color{muted}Focus: .NET, React, Agile, Cloud \& DevOps}

\cventry{sep 2022 -- jun 2024}{Breda, NL}
        {MBO Niveau 4 BBL -- ICT Systeem- \& Devicetechnicus}{Curio}
{\small\color{muted}Versneld duaal traject: fulltime werkzaam bij Dentech gecombineerd met formele ICT-systeemopleiding}
\end{cvgrid}

% ═══════════════════════════════════════════════════════════
% VAARDIGHEDEN, CERTIFICERINGEN & TALEN
% ═══════════════════════════════════════════════════════════
\section{Vaardigheden \& Certificeringen}

\begin{cvgrid}[itemsep=2pt]
\cvrow{Backend}{C\# / .NET, ASP.NET Core, Entity Framework Core, REST API's, PostgreSQL, SQL Server}
\cvrow{Messaging}{RabbitMQ, asynchrone verwerking, event-driven architecturen}
\cvrow{Frontend}{React, TypeScript, JavaScript}
\cvrow{Observability}{OpenTelemetry, metrics, tracing, gestructureerde logging}
\cvrow{Infrastructuur}{Linux, Proxmox, VMware, Hyper-V, Docker, Windows Server, Active Directory, PowerShell, Bash}
\cvrow{Netwerken}{Cisco IOS, VLAN's, inter-VLAN routing, firewalls, VPN, DNS/DHCP, 3CX PBX}
\cvrow{Tools}{Git, CI/CD, Azure, Agile / Scrum}
\vspace{3pt}
\cvrow{Certificeringen}{Microsoft 365 Fundamentals (MS-900) \middot 3CX Advanced Certified Engineer \middot Yealink Junior Certified IP Phone Engineer \middot RoutIT Hosted Voice Fundamentals}
\cvrow{Talen}{Nederlands (moedertaal) \middot Engels (professioneel werkend) \middot Duits (beginnend)}
\end{cvgrid}

\end{document}
